\documentclass[french]{article}


\usepackage{babel}
\usepackage{psfig}
\usepackage{latexsym}
\usepackage{amssymb}
\usepackage{comment}
\usepackage{moreverb}
\usepackage{a4wide}
\usepackage{times}
\usepackage[T1]{fontenc}
\usepackage[french]{minitoc}
\usepackage{latexsym,color}

\newenvironment{solution}{\nopagebreak\par\vspace*{2mm}\underline{Solution~:}\em\par}{\par\vspace*{2mm}}


%\excludecomment{solution}

\begin{document}

\begin{center}
{\Large\bf Examen ENSTA}
\end{center}

%%%%%%%%%%%%%%%%%%%%%%%%%%%%%%%%%%%%%%%%%%%%%%%%%%%%%%%%%%%%

\vspace*{0.5cm}


\section{Indexation}

Soit un fichier de 1~000~000 enregistrements r�partis
        en pages de 4\,000 octets. Chaque enregistrement
        fait 45 octets. R�pondez aux questions suivantes en justifiant
    vos r�ponses (on suppose que les pages sont pleines)\,?
  \begin{enumerate}
   \item Quelle est la taille du fichier\,? Combien faut-il de pages\,? 
\begin{solution}
\begin{enumerate}
  \item 45 Mo!
  \item $\frac{45000000}{4000}=11250$
\end{enumerate}
\end{solution}
   \item Quelle est la taille d'un index de type arbre-B
        si la cl� fait 32 octets
           et un pointeur 8 octets\,?
\begin{solution}
Une entr�e d'index fait 40 octets. Donc on en met
$\frac{4000}{40}=100$ par page. Il faut indexer 1\,000\,000
enregistrements pour le dernier niveau de l'index, soit
$\frac{1000000}{100}=10000$ pages. Ensuite il
faut indexer chacune des 1000 pages, soit $\frac{1000}{100}=10$
pages suppl�mentaires, que pour finir on indexe avec la racine.

Donc il faut 10000 + 10 + 1 pages dans l'index.
\end{solution}
 
   \item  Si on suppose qu'un E/S co�te 10 ms, quel est le co�t moyen
    d'une recherche
     d'un enregistrement par cl� unique,
   avec index et sans index\,?

 \begin{solution}
Recherche avec index\,: trois lectures dans l'index, 
puis une lecture dans le fichier. Soit 40 ms.

Sans index, en moyenne il faut lire la moiti� du fichier,
soit 5625 pages, soit plus de 50 secondes.
\end{solution}
\end{enumerate}

\section{Concurrence}

Le programme suivant s'ex�cute dans un syst�me de gestion de commandes pour
les produits d'une entreprise. Il permet de commander une quantit� donn�e
d'un produit qui se trouve en stock. Les param�tres du programme repr�sentent
respectivement la r�f�rence de la commande ($c$), la r�f�rence du produit ($p$)
et la quantit� command�e ($q$).

\begin{tabbing}
tab \= tab \= \kill
{\bf Commander} ($c$, $p$, $q$)\\
\>  {\em Start}\\
\>  Lecture prix produit $p$\\
\>  \underline{si} $q$ > stock de $p$ \underline{alors} {\em Abort}\\
\>  \underline{sinon}\\
\> \>    Mise � jour stock de $p$\\
\> \>    Enregistrement dans $c$ du total de facturation\\
\> \>    {\em Commit}\\
\>  \underline{fin si}
\end{tabbing}

{\bf N.B.}\,: le prix et la quantit� de produit en 
stock sont gard�s dans des enregistrements diff�rents.

\begin{enumerate}
	\item Lesquelles des transactions suivantes peuvent �tre
obtenues par l'execution du programme ci-dessus? Justifiez votre r�ponse.

{\small
\begin{verbatim}
T: r[x] r[y] a
T': r[x] a
T'': r[x] w[y] w[z] c
T''': r[x] r[y] w[y] w[z] c
\end{verbatim}
}

	\item Dans le syst�me s'ex�cutent en m�me temps trois
transactions: deux commandes d'un m�me produit et le changement du prix
de ce m�me produit. Montrez que l'histoire ci-dessous est une ex�cution
concurrente de ces trois transactions et expliquez la signification des
enregistrements qui y interviennent.

{\bf H} : r$_1$[x] r$_1$[y] w$_2$[x] w$_1$[y] c$_2$ r$_3$[x] r$_3$[y] 
w$_1$[z] c$_1$ w$_3$[y] w$_3$[u] c$_3$

	\item V�rifiez si {\bf H} est s�rialisable en identifiant
les conflits et en construisant le graphe de s�rialisation.

	\item Quelle est l'ex�cution obtenue par verrouillage �
deux phases � partir de {\bf H}? Quel prix sera appliqu� pour la seconde 
commande, le m�me que pour la premi�re ou le prix modifi�?

On consid\`ere que le rel\^achement des verrous d'une transaction se fait au 
{\em Commit} et qu'\`a ce moment on ex\'ecute en priorit\'e les op\'erations 
bloqu\'ees en attente de verrou, dans l'ordre de leur blocage.
\end{enumerate}

\end{document}
